\documentclass[11pt,a4paper]{article}

\usepackage[margin=2.6cm]{geometry}
\usepackage{amsmath,amssymb,amsthm}
\usepackage{booktabs}
\usepackage{microtype}
\usepackage[colorlinks=true,linkcolor=blue,citecolor=blue,urlcolor=blue]{hyperref}

\theoremstyle{plain}

% \newtheorem{theorem}{Theorem}
% \newtheorem{proposition}[theorem]{Proposition}
% \newtheorem{lemma}[theorem]{Lemma}
% \newtheorem{corollary}[theorem]{Corollary}
% \theoremstyle{definition}
% \newtheorem{definition}[theorem]{Definition}
% \newtheorem{remark}[theorem]{Remark}
% \newtheorem{example}{Example}

\newtheorem{theorem}{Theorem}
\newtheorem{proposition}{Proposition}
\newtheorem{lemma}{Lemma}
\newtheorem{corollary}{Corollary}
\theoremstyle{definition}
\newtheorem{definition}{Definition}
\newtheorem{remark}{Remark}
\newtheorem{example}{Example}

\newcommand{\R}{\mathbb{R}}
\newcommand{\SEthree}{\mathrm{SE}(3)}
\newcommand{\rank}{\operatorname{rk}}
\newcommand{\vspan}{\operatorname{span}}
\newcommand{\Bper}{\mathbf{B}^{\times}}
\newcommand{\Bnode}{\mathbf{B}^{\bullet}}
\newcommand{\Ebar}{\bar{\mathbf{E}}}
\newcommand{\Ebaro}{\bar{\mathbf{E}}_o}
\newcommand{\Dp}{\mathbf{D}_p}
\newcommand{\Da}{\mathbf{D}_a}

\title{\bf A Node-indexed Construction of the Extended Bearing Rigidity Matrix for Heterogeneous Networks}
\author{Noyan Erdin Kilic\\
\texttt{noyanerdin.kilic@studenti.unipd.it}}
\date{}

\begin{document}
\maketitle

\begin{abstract}
The unified framework of \cite{michieletto2021} tests infinitesimal bearing
rigidity by comparing the rank of an extended rigidity matrix with the rank of
the same matrix built on the complete graph. While applying it to formations
that mix ground vehicles with aerial platforms, we obtained rank comparisons
that disagreed with rigidity, and this note reports what we believe is the
cause. In eq.~(10) of \cite{michieletto2021} the factor $\mathbf U$ is indexed by edges
and multiplies the incidence matrix from the left, so within a row block it acts on the
difference of the two endpoint variations and imposes the same restriction on
both of them. When the endpoints lie in different manifolds, particularly in one planar and
one spatial, it is not possible to find a choice of $\mathbf U_{ij}$ that both
reproduces the bearing differential and vanishes on the inadmissible
directions of the two endpoints. The two requirements are compatible only when
the agents are vertically aligned. Definition~13 of \cite{michieletto2021} is
still satisfied, since it constrains the matrix only on admissible variations,
but the rank then carries a contribution that depends on the edge set, and it
is ranks that Theorem~2 compares. We give a minimal example checkable by hand, a three-agent formation, that the comparison reports as flexible although it is rigid. Restricting
the degrees of freedom node-wise rather than edge-wise removes the difficulty,
and the resulting matrix coincides with eq.~(10) of \cite{michieletto2021} for
every homogeneous formation, so Table~I and Section~V of
\cite{michieletto2021} are unaffected. The rigidity conclusions of the two
case studies there are also unaffected, since both use the complete graph. A node-wise restriction of this kind is already used in the
heterogeneous formation control formulation of \cite{Pozzan2021}, which
supports the reading proposed here. All statements are checked against a
finite-difference Jacobian of the bearing map.
\end{abstract}

\section{Summary}

Consider three agents at
\begin{equation}\label{eq:ex}
  p_1=(0,0,0),\qquad p_2=(2,0,0),\qquad p_3=(1,1,1),
\end{equation}
with agents $1$ and $2$ in $\mathcal D_1=\mathcal D_2=\R^2$ and
agent $3$ in $\mathcal D_3=\R^3$, and with $R_i = I_3$, the attitudes being constant since neither domain has rotational freedom. Let $\mathcal G$ consist of the two edges
$(v_1,v_3)$ and $(v_2,v_3)$. The formation has $c=7$ degrees of freedom, and
Section~\ref{sec:example} shows that its bearing-preserving motions are exactly
two horizontal translations and uniform scaling, for $\mathcal G$ and for the
complete graph $\mathcal K$. It is therefore infinitesimally bearing rigid.

\begin{center}
\begin{tabular}{lccl}
\toprule
& $\rank\mathbf B_{\mathcal G}$ & $\rank\mathbf B_{\mathcal K}$ & Theorem~2 concludes\\
\midrule
$\Bper_{\mathcal G}$, eq.~(10) of \cite{michieletto2021} & $4$ & $5$ & flexible\\
$\Bnode_{\mathcal G}$, eq.~\eqref{eq:pernode} below        & $4$ & $4$ & rigid\\
finite differences of $b_{\mathcal G}$            & $4$ & $4$ & rigid\\
\bottomrule
\end{tabular}
\end{center}

The extra unit of rank in case of $\Bper_{\mathcal G}$ on the complete graph is associated with a vertical motion of
the two planar agents. A single factor $\mathbf U_{ij}$ sits between $\Dp$ and
$\Ebar^\top$, so it acts on the difference of the two endpoint variations and
applies the same restriction to both endpoints, while a planar and a spatial
agent need different ones. Section~\ref{sec:expand} writes the products out
block by block.

The matrix $\Bper_{\mathcal G}$ reproduces the bearing differential on
admissible variations, so Definition~13 of \cite{michieletto2021} is met, and
the ranks of the two constructions agree once both are restricted to motions
the agents can perform.
The two are identical for every homogeneous formation, and more generally
whenever the two endpoints of each edge have the same translational projector
(Proposition~\ref{prop:homog}), so Table~I and Section~V of
\cite{michieletto2021} are untouched. Both case studies of
\cite{michieletto2021} use complete graphs, where
$\rank\mathbf B_{\mathcal G}=\rank\mathbf B_{\mathcal K}$ holds trivially. The
paper also studies proper subgraphs $\mathcal G_1$ and $\mathcal G_2$ in the
heterogeneous case, but its conclusions about those subgraphs follow from
explicit infinitesimal variations rather than from a rank comparison. Seeing a
discrepancy requires a heterogeneous formation on an incomplete graph (that includes an edge between a planar and spatial agent).

\section{Notation}

We follow \cite{michieletto2021}. Agent $i$ has position $p_i$, attitude $R_i$
and $c_i=\dim\mathcal D_i$ degrees of freedom, and $c=\sum_ic_i$. An edge
$e_k=(v_i,v_j)$ carries the bearing measured by $i$ of $j$,
\begin{equation}\label{eq:bearing}
  b_{ij}=R_i^\top\hat p_{ij},
  \qquad
  \hat p_{ij}=\frac{p_j-p_i}{\lVert p_j-p_i\rVert},
  \qquad
  s_{ij}=\lVert p_j-p_i\rVert^{-1}.
\end{equation}

We differentiate \eqref{eq:bearing} in the two factors separately, since the
two carry the two halves of the rigidity matrix. For the direction, with
$p_{ij}=p_j-p_i$,
\begin{align}
  \frac{d}{dt}\,\hat p_{ij}
  &=\frac{\dot p_{ij}}{\lVert p_{ij}\rVert}
    -\frac{p_{ij}\,\bigl(p_{ij}^\top\dot p_{ij}\bigr)}{\lVert p_{ij}\rVert^{3}}
   =s_{ij}\Bigl(I_3-\hat p_{ij}\hat p_{ij}^\top\Bigr)\dot p_{ij}
   \notag\\[2pt]
  &=s_{ij}P(\hat p_{ij})\,(\dot p_j-\dot p_i),
  \qquad P(x)=I_3-xx^\top . \label{eq:dphat}
\end{align}
For the attitude, with $\omega_i$ the angular velocity in the world frame we
have $\dot R_i=[\omega_i]_\times R_i$, hence
$\dot R_i^\top=R_i^\top[\omega_i]_\times^\top=-R_i^\top[\omega_i]_\times$ and
\begin{equation}\label{eq:dR}
  \dot R_i^\top\hat p_{ij}
  =-R_i^\top[\omega_i]_\times\hat p_{ij}
  =-R_i^\top\bigl(\omega_i\times\hat p_{ij}\bigr)
  =+R_i^\top\bigl(\hat p_{ij}\times\omega_i\bigr)
  =R_i^\top[\hat p_{ij}]_\times\,\omega_i .
\end{equation}
Adding \eqref{eq:dphat} and \eqref{eq:dR},
\begin{equation}\label{eq:diff}
  \dot b_{ij}
  =\underbrace{s_{ij}R_i^\top P(\hat p_{ij})}_{\textstyle\Dp^{(k)}}
     (\dot p_j-\dot p_i)
   \;-\;\underbrace{\bigl(-R_i^\top[\hat p_{ij}]_\times\bigr)}_{\textstyle\Da^{(k)}}\,\omega_i ,
\end{equation}
which is eq.~(1) of \cite{michieletto2021} with the sign of $\Da^{(k)}$ as
defined there. The minus in front of $\Da^{(k)}$ is the one supplied by
$\Ebaro^\top$ in Section~\ref{sec:expand}.

Let $S_i$ and $P_i$ be the orthogonal projectors onto the translational and
rotational directions agent $i$ can realise, so that
$\rank S_i+\rank P_i=c_i$.

\begin{center}
\begin{tabular}{lccccc}
\toprule
$\mathcal D_i$ & $\R^2$ & $\R^3$ & $\R^2\times S^1$ & $\R^3\times S^1$ & $\SEthree$\\
\midrule
$S_i$ & $\mathrm{diag}(1,1,0)$ & $I_3$ & $\mathrm{diag}(1,1,0)$ & $I_3$ & $I_3$\\
$P_i$ & $0$ & $0$ & $e_3e_3^\top$ & $vv^\top$ & $I_3$\\
\bottomrule
\end{tabular}
\end{center}

A variation $\delta$ is lifted to $\delta^+\in\R^{6n}$ by zero padding, as in
(8) of \cite{michieletto2021}. The lifted variations fill the
\emph{admissible subspace}
\begin{equation}\label{eq:A}
  \mathcal A=\operatorname{range}\bar S\oplus\operatorname{range}\bar P
  \subseteq\R^{6n},
  \qquad
  \bar S=\mathrm{blkdiag}(S_i),\quad \bar P=\mathrm{blkdiag}(P_i),
\end{equation}
of dimension $c$. Let $A\in\R^{6n\times c}$ have orthonormal columns spanning
$\mathcal A$, and take the variation coordinates to be the ones it induces, so
that
\begin{equation}\label{eq:embed}
  \delta^+=A\delta .
\end{equation}
We write $q_v:=6n-c$ for the dimension of $\mathcal A^\perp$. Finally let
$J_{\mathcal G}\in\R^{3m\times c}$ denote the Jacobian of the bearing map in
these coordinates, so that $J_{\mathcal G}\delta=\dot b_{\mathcal G}$ whenever
$\delta^+=A\delta$, with $\dot b_{\mathcal G}$ stacking the row blocks
\eqref{eq:diff}.

\section{Where the DOF restriction enters the product}\label{sec:expand}

We write $\Bper_{\mathcal G}$ for the extended matrix of Proposition~2 of
\cite{michieletto2021} and $\Bnode_{\mathcal G}$ for the node-indexed
construction, and use these names throughout. The first is
\begin{equation}\label{eq:peredge}
  \Bper_{\mathcal G}=\bigl[\;\Dp\,\mathbf U\,\Ebar^\top\quad
                             \Da\,\mathbf V\,\Ebaro^\top\;\bigr],
  \qquad
  \mathbf U=\mathrm{diag}(\mathbf U_{ij}),\quad
  \mathbf V=\mathrm{diag}(\mathbf V_{ij}),
\end{equation}
and the second is the same object with the two factors moved to the right of
the incidence matrices and indexed by node,
\begin{equation}\label{eq:pernode}
  \Bnode_{\mathcal G}=\bigl[\;\Dp\,\Ebar^\top\,\bar S\quad
                              \Da\,\Ebaro^\top\,\bar P\;\bigr].
\end{equation}

\subsection*{The incidence matrices}

For a directed graph, \cite{michieletto2021} defines the incidence matrix
$\mathbf E\in\R^{n\times m}$ and the matrix $\mathbf E_o\in\R^{n\times m}$
collecting its outgoing part by
\begin{equation}\label{eq:inc}
\begin{aligned}
  [\mathbf E]_{ik}&=
  \begin{cases}
    -1,&\text{if }e_k=(v_i,v_j)\in\mathcal E\ \text{(outgoing edge)}\\
    \ \ 1,&\text{if }e_k=(v_j,v_i)\in\mathcal E\ \text{(ingoing edge)}\\
    \ \ 0,&\text{otherwise,}
  \end{cases}\\[6pt]
  [\mathbf E_o]_{ik}&=
  \begin{cases}
    -1,&\text{if }e_k=(v_i,v_j)\in\mathcal E\ \text{(outgoing edge)}\\
    \ \ 0,&\text{otherwise,}
  \end{cases}
\end{aligned}
\end{equation}
which are eqs.~(1) and~(2) of \cite{michieletto2021}, together with
$\Ebar=\mathbf E\otimes I_3$ and $\Ebaro=\mathbf E_o\otimes I_3$. Transposing
turns the $k$th column into the $k$th row block, so for an edge
$e_k=(v_i,v_j)$, carrying the $-1$ at its tail $v_i$ and the $+1$ at its head
$v_j$,
\begin{align}
  \bigl[\Ebar^\top\bigr]_{k,\,\cdot}
  &=\Bigl[\ \mathbf 0\ \cdots\ \underbrace{-I_3}_{\text{agent }i}\ \cdots\
     \underbrace{+I_3}_{\text{agent }j}\ \cdots\ \mathbf 0\ \Bigr],
  \label{eq:Erow}\\[4pt]
  \bigl[\Ebaro^\top\bigr]_{k,\,\cdot}
  &=\Bigl[\ \mathbf 0\ \cdots\ \underbrace{-I_3}_{\text{agent }i}\ \cdots\
     \underbrace{\ \ \mathbf 0\ \ }_{\text{agent }j}\ \cdots\ \mathbf 0\ \Bigr].
  \label{eq:Eorow}
\end{align}
A bearing depends on the positions of \emph{both} endpoints, so $\Ebar^\top$
touches both column blocks, whereas it depends on the attitude of the sensing
agent only, so $\Ebaro^\top$ touches one.

\subsection*{The edge-indexed product}

$\Dp$ and $\mathbf U$ are block diagonal over \emph{edges}, so multiplying
\eqref{eq:Erow} from the left by $\Dp^{(k)}\mathbf U_{ij}$ scales the whole row
block,
\begin{equation}\label{eq:perrow}
  \bigl[\Dp\,\mathbf U\,\Ebar^\top\bigr]_{k,\,\cdot}
  =\Bigl[\ \mathbf 0\ \cdots\
     \underbrace{-\Dp^{(k)}\mathbf U_{ij}}_{\text{agent }i}\ \cdots\
     \underbrace{+\Dp^{(k)}\mathbf U_{ij}}_{\text{agent }j}\ \cdots\
     \mathbf 0\ \Bigr],
\end{equation}
and likewise, using \eqref{eq:Eorow},
\begin{equation}\label{eq:perrowa}
  \bigl[\Da\,\mathbf V\,\Ebaro^\top\bigr]_{k,\,\cdot}
  =\Bigl[\ \mathbf 0\ \cdots\
     \underbrace{-\Da^{(k)}\mathbf V_{ij}}_{\text{agent }i}\ \cdots\
     \underbrace{\ \ \mathbf 0\ \ }_{\text{agent }j}\ \cdots\ \mathbf 0\ \Bigr].
\end{equation}
Acting on $\delta^+$ this returns
\begin{equation}\label{eq:actper}
  (\Bper\delta^+)_k
  =\Dp^{(k)}\mathbf U_{ij}\,\delta_{p,j}-\Dp^{(k)}\mathbf U_{ij}\,\delta_{p,i}
   -\Da^{(k)}\mathbf V_{ij}\,\delta_{a,i}
  =\Dp^{(k)}\mathbf U_{ij}\bigl(\delta_{p,j}-\delta_{p,i}\bigr)
   -\Da^{(k)}\mathbf V_{ij}\,\delta_{a,i},
\end{equation}
which agrees with \eqref{eq:diff} when $\mathbf U_{ij}$ acts as the identity on
the relative motion. Comparing the last two expressions in \eqref{eq:actper},
the same matrix $\Dp^{(k)}\mathbf U_{ij}$ multiplies $\delta_{p,j}$ and
$\delta_{p,i}$, and it may equivalently be read as acting on their difference.
That is the whole of Lemma~\ref{lem:same} below.

Note also from \eqref{eq:perrowa} that $\mathbf V_{ij}$ meets only agent $i$.
Setting $\mathbf V_{ij}=P_i$ restricts only the agent whose columns it occupies
and does not cause problems. Everything below concerns $\mathbf U$ alone.

\subsection*{The node-indexed product}

Now multiply \eqref{eq:Erow} from the \emph{right} by
$\bar S=\mathrm{blkdiag}(S_1,\dots,S_n)$, which is block diagonal over
\emph{nodes}, so each column block is scaled by its own projector,
\begin{equation}\label{eq:noderow}
  \bigl[\Ebar^\top\bar S\bigr]_{k,\,\cdot}
  =\Bigl[\ \mathbf 0\ \cdots\ -S_i\ \cdots\ +S_j\ \cdots\ \mathbf 0\ \Bigr],
  \qquad
  \bigl[\Dp\,\Ebar^\top\bar S\bigr]_{k,\,\cdot}
  =\Bigl[\ \mathbf 0\ \cdots\ -\Dp^{(k)}S_i\ \cdots\ +\Dp^{(k)}S_j\ \cdots\ \mathbf 0\ \Bigr],
\end{equation}
and in the same way
$\bigl[\Da\Ebaro^\top\bar P\bigr]_{k,\,\cdot}$ has the single nonzero block
$-\Da^{(k)}P_i$ in the columns of agent $i$. Hence
\begin{equation}\label{eq:actnode}
  (\Bnode\delta^+)_k
  =\Dp^{(k)}\bigl(S_j\delta_{p,j}-S_i\delta_{p,i}\bigr)
   -\Da^{(k)}P_i\,\delta_{a,i}.
\end{equation}
Comparing \eqref{eq:actper} with \eqref{eq:actnode}, the difference is only
whether the projector is applied before or after the two endpoint variations
are subtracted. If $S_i=S_j=S$, then
$S_j\delta_{p,j}-S_i\delta_{p,i}=S(\delta_{p,j}-\delta_{p,i})$ and the two row
blocks agree. This condition is weaker than homogeneity, since only the
translational projector enters it. An edge joining an $\R^3\times S^1$ agent to
an $\SEthree$ one has $S_i=S_j=I_3$, and an edge joining an $\R^2$ agent to an
$\R^2\times S^1$ one has $S_i=S_j=\mathrm{diag}(1,1,0)$, so both are covered.
Proposition~\ref{prop:homog} states the general case.

\subsection*{The two requirements}

\begin{lemma}[the two endpoints share one factor]\label{lem:same}
By \eqref{eq:perrow}, the column blocks of agents $i$ and $j$ inside row block
$k$ of \eqref{eq:peredge} are $-\Dp^{(k)}\mathbf U_{ij}$ and
$+\Dp^{(k)}\mathbf U_{ij}$. Any restriction encoded by $\mathbf U_{ij}$ is
therefore imposed on both endpoints at once.
\end{lemma}

Two requirements are placed on the column blocks of the bearing rigidity
matrix. The matrix should reproduce \eqref{eq:diff} on motions the agents can
perform, and the text preceding Theorem~2 in \cite{michieletto2021} asks that
inadmissible degrees of freedom index null columns of the matrix (so that the
count $q_v$ is a property of the manifolds alone and cancels between
$\rank\Bper_{\mathcal G}=6n-q_v-q_i$ and
$\rank\Bper_{\mathcal K}=6n-q_v-q_t$). Written for a candidate matrix
$\mathbf B\in\R^{3m\times6n}$, the two requirements read
\begin{align}
  \mathbf BA&=J_{\mathcal G}, \tag{F}\label{eq:F}\\
  \mathbf B(I-AA^\top)&=0. \tag{Z}\label{eq:Z}
\end{align}

Here $q_i$ and $q_t$ denote the dimensions of the infinitesimal and the trivial
variation sets of \cite{michieletto2021}, both read in the coordinates of
$\delta$, so that $q_i=\dim\ker(\mathbf B_{\mathcal G}A)$ and
$q_t=\dim\ker(\mathbf B_{\mathcal K}A)$.

Definition~13 of \cite{michieletto2021} determines a matrix only on
$\mathcal A$, since $\delta^+$ never leaves that subspace, so it does not
single out either construction. Both $\Bper_{\mathcal G}$ and
$\Bnode_{\mathcal G}$ satisfy it, and they differ only on $\mathcal A^\perp$.
Condition \eqref{eq:F} is Definition~13 restated, and \eqref{eq:Z} is what
decides between them.


\begin{lemma}[when \eqref{eq:F} and \eqref{eq:Z} are compatible]\label{lem:obstruction}
Let $e_k=(v_i,v_j)$ with $i$ planar ($S_i=\mathrm{diag}(1,1,0)$) and $j$
spatial ($S_j=I_3$). Then $\Bper_{\mathcal G}$ satisfies \eqref{eq:F} and
\eqref{eq:Z} on the column blocks of $i$ and $j$ only if $P(\hat p_{ij})e_3=0$,
that is, only if agent $j$ lies vertically above or below agent $i$.
\end{lemma}

\begin{proof}
Write $U=\mathbf U_{ij}$ and read the two column blocks of \eqref{eq:perrow} in
turn.

\emph{Agent $j$.} Its admissible directions are all of $\R^3$, and by
\eqref{eq:actnode} the block that reproduces \eqref{eq:diff} there is
$+\Dp^{(k)}S_j=\Dp^{(k)}$. So \eqref{eq:F} reads
\begin{equation}\label{eq:step1}
  \Dp^{(k)}U\,x=\Dp^{(k)}x\quad\text{for all }x\in\R^3,
  \qquad\text{i.e.}\qquad
  \Dp^{(k)}U=\Dp^{(k)} .
\end{equation}

\emph{Agent $i$.} Its inadmissible direction is $e_3$, and by \eqref{eq:perrow}
the corresponding column of the matrix is $-\Dp^{(k)}Ue_3$. So \eqref{eq:Z}
reads
\begin{equation}\label{eq:step2}
  \Dp^{(k)}U\,e_3=0 .
\end{equation}

Substituting \eqref{eq:step1} into \eqref{eq:step2} gives
$\Dp^{(k)}e_3=0$, that is $e_3\in\ker\Dp^{(k)}$. Finally
\begin{equation*}
  \Dp^{(k)}=s_{ij}R_i^\top P(\hat p_{ij}),
  \qquad s_{ij}\neq0,\quad R_i\ \text{orthogonal},
  \qquad\text{so}\qquad
  \ker\Dp^{(k)}=\ker P(\hat p_{ij})=\vspan\{\hat p_{ij}\},
\end{equation*}
and $e_3\in\vspan\{\hat p_{ij}\}$ says $e_3=c\,\hat p_{ij}$ for some $c\in\R$.
Both vectors have unit norm, so $c=\pm1$ and $\hat p_{ij}=\pm e_3$.
\end{proof}

Running the same two steps with general $S_i,S_j$ gives
$\Dp^{(k)}U=\Dp^{(k)}S_i$ from agent $i$ and $\Dp^{(k)}U=\Dp^{(k)}S_j$ from
agent $j$, hence
\begin{equation}\label{eq:general}
  \Dp^{(k)}(S_i-S_j)=0,
  \qquad\text{i.e.}\qquad
  \operatorname{range}(S_i-S_j)\subseteq\vspan\{\hat p_{ij}\},
\end{equation}
a nontrivial algebraic condition on $(p_i,p_j)$ whenever $S_i\neq S_j$.

\begin{remark}[why this stays invisible until two graphs are compared]
By \eqref{eq:F} both constructions have the same rank once restricted to
admissible columns, $\rank(\Bper A)=\rank(\Bnode A)=\rank J$. They differ only
in how much rank the inadmissible columns add, and that amount depends on the
edge set. Theorem~2 compares $\rank\mathbf B_{\mathcal G}$ with
$\rank\mathbf B_{\mathcal K}$, so a difference in that amount is read as a
difference in rigidity.
\end{remark}

\section{Analysis of the example network}\label{sec:example}

The formation \eqref{eq:ex} has $c=7$ degrees of freedom, two horizontal
components for each of the planar agents and three for agent~3. Write a
variation as $\delta=(u_1,u_2,u_3)$ with $(u_1)_3=(u_2)_3=0$ and
$u_3\in\R^3$.

No agent rotates and $R_i=I_3$, so \eqref{eq:diff} reduces to
$\dot b_{ij}=s_{ij}P(\hat p_{ij})(u_j-u_i)$, and $\ker P(x)=\vspan\{x\}$. A
bearing is preserved to first order exactly when its two endpoints separate
along the line joining them. For the two edges of $\mathcal G$,
\begin{equation}\label{eq:exker}
  u_3-u_1=\alpha\,(p_3-p_1)=\alpha\,(1,1,1),
  \qquad
  u_3-u_2=\beta\,(p_3-p_2)=\beta\,(-1,1,1),
  \qquad \alpha,\beta\in\R ,
\end{equation}
where $\alpha$ and $\beta$ are rates, the spanning vector being taken
unnormalised so that the two third components are both $1$.

\emph{Third components.} Both $p_3-p_1$ and $p_3-p_2$ have third component $1$,
while $u_1$ and $u_2$ have none. Reading the third row of each equation in
\eqref{eq:exker} therefore gives $\alpha=(u_3)_3$ and $\beta=(u_3)_3$, so the
two scalars are equal. Call the common value $t$.

\emph{Both equations at once.} With $\alpha=\beta=t$, \eqref{eq:exker} reads
$u_i=u_3-t\,(p_3-p_i)$ for $i=1,2$, that is
\begin{equation}\label{eq:exsim}
  u_i=w+t\,p_i \quad (i=1,2,3),
  \qquad w:=u_3-t\,p_3 ,
\end{equation}
which describes an infinitesimal similarity, translating the formation by $w$
and scaling it by $t$ about the origin.

\emph{Which similarities are admissible.} Since $p_1=0$ we get $w=u_1$, which
is horizontal. Conversely any $(w,t)$ with $w$ horizontal is admissible,
because $u_2=w+t\,p_2$ is horizontal, $p_2$ being horizontal, and $u_3$ is
unconstrained. The kernel is thus $3$-dimensional, spanned by the two
horizontal translations and the uniform scaling.

\emph{The same for $\mathcal K$.} A similarity \eqref{eq:exsim} preserves
\emph{every} bearing of the formation, since $u_j-u_i=t\,(p_j-p_i)$ for all
$i,j$, not only for the two edges of $\mathcal G$. Hence
$\ker J_{\mathcal G}=\ker J_{\mathcal K}$, the framework is IBR in the sense of
Definition~12 of \cite{michieletto2021}, and
\begin{equation*}
  \rank J_{\mathcal G}=\rank J_{\mathcal K}=c-3=4 .
\end{equation*}
Two edges, each contributing rank $2$, attain this, so $\mathcal G$ is minimally
rigid.

For $e_1=(v_1,v_3)$ we have $p_3-p_1=(1,1,1)$ of norm $\sqrt3$ and
$s_{13}=1/\sqrt3$, so
\begin{equation*}
  \Dp^{(1)}=s_{13}\bigl(I_3-\hat p_{13}\hat p_{13}^\top\bigr)
  =\frac{1}{\sqrt3}\left(I_3-\frac13
  \begin{bmatrix}1&1&1\\1&1&1\\1&1&1\end{bmatrix}\right)
  =\frac{\sqrt3}{9}
  \begin{bmatrix}\ 2&-1&-1\\-1&\ 2&-1\\-1&-1&\ 2\end{bmatrix},
\end{equation*}
and for $e_2=(v_2,v_3)$ we have $p_3-p_2=(-1,1,1)$, again of norm $\sqrt3$, so
\begin{equation*}
  \Dp^{(2)}
  =\frac{\sqrt3}{9}
  \begin{bmatrix}2&1&1\\1&2&-1\\1&-1&2\end{bmatrix}.
\end{equation*}
Placing these by \eqref{eq:perrow} with $\mathbf U_{ij}=I_3$, and dropping the
attitude half which is identically zero here,
\begin{equation}\label{eq:exmat}
  \Bper_{\mathcal G}
  =\frac{\sqrt3}{9}
  \left[\begin{array}{rrr|rrr|rrr}
    -2& 1& 1& 0& 0& 0&  2&-1&-1\\
     1&-2& 1& 0& 0& 0& -1& 2&-1\\
     1& 1&-2& 0& 0& 0& -1&-1& 2\\\hline
     0& 0& 0&-2&-1&-1&  2& 1& 1\\
     0& 0& 0&-1&-2& 1&  1& 2&-1\\
     0& 0& 0&-1& 1&-2&  1&-1& 2
  \end{array}\right],
\end{equation}
the columns being $(u_{1x},u_{1y},u_{1z}\mid u_{2x},u_{2y},u_{2z}\mid
u_{3x},u_{3y},u_{3z})$. The third and sixth columns hold the vertical motions
of the two planar agents, and neither is zero. The node-indexed matrix
\eqref{eq:noderow} is \eqref{eq:exmat} with those two columns replaced by
zeros, since $S_1=S_2=\mathrm{diag}(1,1,0)$ annihilates them while $S_3=I_3$
leaves the rest alone.

On $\mathcal G$ both constructions have rank $4$, as two row blocks of rank $2$
cannot give more. The complete graph adds a row block for $e_3=(v_1,v_2)$, with
$\hat p_{12}=e_1$ and $s_{12}=\tfrac12$.\footnote{Since all three agents here
are particles with $R_i=I_3$, reversing an edge negates its row block and
changes no rank or kernel, so $\mathcal K$ may be read as oriented with three
row blocks or as directed with six. Section~V-A1 of \cite{michieletto2021}
takes the graph in the $\R^d$ case to be undirected for the same reason.}

Here the two constructions part. On the nine lifted columns the kernel of
$\Bper_{\mathcal K}$ consists of the four infinitesimal similarities of $\R^3$,
three translations and the scaling, so $\rank\Bper_{\mathcal K}=9-4=5$. Only
three of the four are admissible, since the vertical translation lifts the
planar agents out of their plane. The node-indexed matrix has zeros in columns
three and six, so its rank is that of the remaining seven columns,
$\rank\Bnode_{\mathcal K}=7-3=4$. The comparison of $\Bper_{\mathcal G}$ with $\Bper_{\mathcal K}$ is then $4$
against $5$ and reports flexibility, where that of $\Bnode_{\mathcal G}$ with
$\Bnode_{\mathcal K}$ is $4$ against $4$.

One might suspect that the value $5$ comes from having used
$\mathbf U_{12}=I_3$ on the edge between the two planar agents, where
Table~III of \cite{michieletto2021} prescribes $[\,e_1\ e_2\ 0_3\,]$ for a pair
of terrestrial agents. It does not. The $z$ columns of agents $1$ and $2$
receive entries from three row blocks, one for each edge of $\mathcal K$, and
$\mathbf U_{12}$ governs only one of them. Taking
$\mathbf U_{12}=[\,e_1\ e_2\ 0_3\,]$ turns the factor
$\Dp^{(3)}\mathbf U_{12}$ of that row block from
$\tfrac12\mathrm{diag}(0,1,1)$ into $\tfrac12\mathrm{diag}(0,1,0)$, so it
contributes nothing to the two columns. The row blocks of the two mixed edges
still do, since Table~III prescribes $\mathbf U_{ij}=I_3$ there, and
$\rank\Bper_{\mathcal K}$ is again $5$.

\begin{remark}[a second mechanism]
In the example, the number of null columns happens to be the same for
$\mathcal G$ and $\mathcal K$, and the discrepancy is entirely in how much rank
the nonzero inadmissible columns contribute. The count itself can also become
graph-dependent. For $n=4$ with $\mathcal D_{1,2,3}=\R^2$,
$\mathcal D_4=\R^3$ at $p_1=(0,0,0)$, $p_2=(1,0,0)$, $p_3=(1,1,0)$,
$p_4=(\tfrac12,\tfrac12,1)$ and $\mathcal E=\{(1,2),(1,3),(2,4),(3,4)\}$, every
edge of $\mathcal G$ incident to agent $1$ joins it to another planar agent, so
by Proposition~\ref{prop:homog} those row blocks agree with the node-indexed
ones, where $S_1=\mathrm{diag}(1,1,0)$ sets the $z$ column of agent $1$ to
zero. In $\mathcal K$ that agent is joined to the spatial one as well, its $z$
column is no longer zero, and there are $13$ null columns against $12$. That framework is again minimally rigid ($c-q_t=9-3=6$,
attained by four edges) and again reported flexible,
$\rank\Bper_{\mathcal G}=6$ against $\rank\Bper_{\mathcal K}=8$. Deleting the
three planar $z$ columns of $\Bper_{\mathcal K}$ returns its rank to $6$.
\end{remark}

\begin{remark}[the same arithmetic in Section~VI-B of \cite{michieletto2021}]
The heterogeneous case study there reports $q_t=5$, while the accompanying list
of trivial motions has four entries. Two planar translations, one coordinated
rotation, one scaling. The fifth is the global $z$ translation, which is not admissible
for the planar unicycles. Using the entries of Table~III of
\cite{michieletto2021}, and generic positions since the article does not give
the configuration, we reproduce the reported
$\rank\Bper_{\mathcal K}=13$ with six null columns. Deleting the three
unicycle $z$ columns lowers it to $11=c-q_t=15-4$. The graph there is complete,
so the rigidity conclusion of the case study is unaffected either way.
\end{remark}

\section{Analysis of the node-indexed construction}

\begin{proposition}[uniqueness]\label{prop:unique}
There is exactly one matrix satisfying \eqref{eq:F} and \eqref{eq:Z}, namely
$J_{\mathcal G}A^\top$, and it equals $\Bnode_{\mathcal G}$ of
\eqref{eq:pernode}.
\end{proposition}

\begin{proof}
Since $\bar S$ and $\bar P$ are orthogonal projectors and $\mathcal A$ splits
along the position and the attitude coordinate blocks, the orthogonal projector
onto $\mathcal A$ is
\begin{equation}\label{eq:AAt}
  AA^\top=\mathrm{blkdiag}(\bar S,\bar P).
\end{equation}
Write $\mathbf M:=\bigl[\;\Dp\,\Ebar^\top\quad\Da\,\Ebaro^\top\;\bigr]$ for the
differential before any restriction, so that \eqref{eq:pernode} reads
$\Bnode_{\mathcal G}=\mathbf M\,\mathrm{blkdiag}(\bar S,\bar P)=\mathbf M
AA^\top$ by \eqref{eq:AAt}. Idempotency of $AA^\top$ and $A^\top A=I_c$ then
give
\begin{equation*}
  \Bnode_{\mathcal G}\bigl(I-AA^\top\bigr)
  =\mathbf M\,AA^\top\bigl(I-AA^\top\bigr)=0,
  \qquad
  \Bnode_{\mathcal G}A=\mathbf M\,AA^\top A=\mathbf M A=J_{\mathcal G},
\end{equation*}
the last equality because $\mathbf M$ reproduces \eqref{eq:diff} on admissible
variations. So $\Bnode_{\mathcal G}$ satisfies \eqref{eq:F} and \eqref{eq:Z}.
Conversely, let $\mathbf B$ satisfy both. From
$I_{6n}=AA^\top+(I_{6n}-AA^\top)$ we obtain
\begin{equation*}
  \mathbf B=\mathbf BAA^\top+\mathbf B\bigl(I_{6n}-AA^\top\bigr)
  =\bigl(\mathbf BA\bigr)A^\top=J_{\mathcal G}A^\top ,
\end{equation*}
so the matrix is unique and equals $\Bnode_{\mathcal G}$.
\end{proof}

The construction $\Bnode_{\mathcal G}$ completes Definition~13 of
\cite{michieletto2021} to the virtual subspace that the null-column requirement
of Theorem~2 asks for, and by Proposition~\ref{prop:unique} it is the only one.

\begin{theorem}\label{thm:repair}
For every $\mathcal G$,
$\rank\Bnode_{\mathcal G}=6n-q_v-q_i$ with $q_v=6n-c$ and
$q_i=\dim\ker(\Bnode_{\mathcal G}A)$. Hence
$\rank\Bnode_{\mathcal G}=\rank\Bnode_{\mathcal K}$ if and only if $q_i=q_t$,
that is, if and only if $(\mathcal G,\chi)$ is IBR.
\end{theorem}

\begin{proof}
By \eqref{eq:Z}, $\Bnode_{\mathcal G}=\Bnode_{\mathcal G}AA^\top$, so the row
space is unchanged by restricting to the admissible columns and
\begin{equation*}
  \rank\Bnode_{\mathcal G}
  =\rank\bigl(\Bnode_{\mathcal G}A\bigr)
  =\underbrace{c}_{\text{columns of }A}-\ \underbrace{q_i}_{\dim\ker}
  =(6n-q_v)-q_i .
\end{equation*}
The same holds for $\mathcal K$ with $q_t=\dim\ker(\Bnode_{\mathcal K}A)$ and
with the \emph{same} $q_v$, since by \eqref{eq:A} it depends only on the
manifolds. Subtracting,
$\rank\Bnode_{\mathcal G}-\rank\Bnode_{\mathcal K}=q_t-q_i$, so the two ranks
agree if and only if $q_i=q_t$. By Theorem~1 of \cite{michieletto2021},
$\ker\mathbf B_{\mathcal K}\subseteq\ker\mathbf B_{\mathcal G}$, so equality of
the two dimensions is equality of the two kernels, which is Definition~12 of
\cite{michieletto2021}.
\end{proof}

\begin{proposition}[when the two constructions coincide]\label{prop:homog}
Take $\mathbf U_{ij}=S_i$ and $\mathbf V_{ij}=P_i$. For an edge
$e_k=(v_i,v_j)$ with $S_i=S_j$, the $k$th row blocks of $\Bper_{\mathcal G}$
and $\Bnode_{\mathcal G}$ agree. Hence if $S_i=S_j$ on every edge of
$\mathcal E$, then $\Bper_{\mathcal G}=\Bnode_{\mathcal G}$ at every
configuration. Homogeneous formations are the special case
$\mathcal D_i=\mathcal D$ for all $i$, where $\mathbf U_{ij}=S$ and
$\mathbf V_{ij}=P$ are the entries of Table~I of \cite{michieletto2021}.
\end{proposition}

\begin{proof}
Compare \eqref{eq:actper} and \eqref{eq:actnode} on row block $k$. Writing $S$
for the common value of $S_i$ and $S_j$,
\begin{equation*}
  \Dp^{(k)}\,S\,\bigl(\delta_{p,j}-\delta_{p,i}\bigr)
  =\Dp^{(k)}\bigl(S\delta_{p,j}-S\delta_{p,i}\bigr)
  =\Dp^{(k)}\bigl(S_j\delta_{p,j}-S_i\delta_{p,i}\bigr)
\end{equation*}
by linearity of $S$, and the attitude blocks agree already, since
$\Ebaro^\top$ meets only agent $i$ and $\mathbf V_{ij}=P_i$. Applying this to
every $k$ gives the second claim. In the homogeneous case
$\bar S=I_n\otimes S$ commutes past the incidence matrix,
$\mathbf U\Ebar^\top=\Ebar^\top\bar S$ and
$\mathbf V\Ebaro^\top=\Ebaro^\top\bar P$.\footnote{For
$\R^3\times S^1$ read $\mathbf V_{ij}$ in the projector form $P_i=vv^\top$
rather than the $[\,0_{3\times2}\ v\,]$ form of Table~I of
\cite{michieletto2021}. The two differ as matrices but not in rank, nor in
$q_v$. Only the phrase ``null column'' is sensitive to the choice, which is why
\eqref{eq:Z} is stated as an identity.}
\end{proof}

\begin{remark}[what the condition asks of $\mathcal G$ and of $\mathcal K$]
The hypothesis of Proposition~\ref{prop:homog} is about the edges present in
$\mathcal G$, which is where the graph dependence enters. Since $\mathcal K$
contains every pair, $\Bper_{\mathcal K}=\Bnode_{\mathcal K}$ requires
$S_i=S_j$ for all $i$ and $j$, so any formation mixing planar with spatial
agents fails it on the complete graph while $\mathcal G$ itself may satisfy it.
The converse implication does not hold at the level of ranks. In
Section~\ref{sec:example} both edges of $\mathcal G$ join a planar to a spatial
agent, so $\Bper_{\mathcal G}\neq\Bnode_{\mathcal G}$, and the two still have
the same rank $4$.
\end{remark}

\paragraph{Implementation.}
By \eqref{eq:noderow} the change is one line per block. For row block
$k=(i,j)$, each index naming the top-left corner of a $3\times3$ block,
\begin{equation*}
  \Bnode[\,3k,\,3j\,]\mathrel{+}=\Dp^{(k)}S_j,
  \qquad
  \Bnode[\,3k,\,3i\,]\mathrel{-}=\Dp^{(k)}S_i,
  \qquad
  \Bnode[\,3k,\,3n{+}3i\,]\mathrel{+}=R_i^\top[\hat p_{ij}]_\times P_i ,
\end{equation*}
the last line being $-\Da^{(k)}P_i$ written out via \eqref{eq:diff}. The two
dense $3m\times3m$ factors are no longer needed. One may also avoid the lift
altogether and apply the rank test to the $c$-column matrices
$J_{\mathcal G}=\Bnode_{\mathcal G}A$ and $J_{\mathcal K}$, in which case
$q_v$ never enters the bookkeeping at all.

\paragraph{Relation to the heterogeneous formulation of Pozzan et al.}
The node-wise treatment is not new. It is used in \cite{Pozzan2021}, where each
agent carries its own selection matrices $S_{p,i}\in\R^{3\times c_{p,i}}$ and
$S_{o,i}\in\R^{3\times c_{o,i}}$ mapping its command vector to its linear and
angular velocity, assembled block-diagonally into
$\mathbf S_p=\mathrm{diag}(S_{p,i})$ and $\mathbf S_o=\mathrm{diag}(S_{o,i})$
in eqs.~(3), (4) and (7) there. The bearing rigidity matrix defined there is
$\mathbf B_{\mathcal G}(\mathbf x)\in\R^{3m\times c}$, and its $k$th row block,
eq.~(9) of \cite{Pozzan2021}, carries $S_{p,i}$ in the column block of agent
$i$ and $S_{p,j}$ in the column block of agent $j$. The two endpoints are
restricted separately, which is what \eqref{eq:noderow} does here.
Infinitesimal rigidity is then defined by
$\ker\mathbf B_{\mathcal K}=\ker\mathbf B_{\mathcal G}$ on those $c$ columns,
so neither the lift nor the count $q_v$ ever appears.

The correspondence with \eqref{eq:pernode} is
$J_{\mathcal G}=\Bnode_{\mathcal G}A$, with
$\mathrm{blkdiag}(\mathbf S_p,\mathbf S_o)$ in the role of $A$. Note that
$S_{p,i}$ is an injection whose columns span the admissible directions rather
than a projector, so it plays the role of $A$ and not of $\bar S$, the relation
being $S_i=S_{p,i}S_{p,i}^\top$ for orthonormal columns. The two formulations
differ also in that \cite{Pozzan2021} parametrises the commands in the body
frame, $\dot p_i=R_iv_i$, which is what puts $R_i^\top R_jS_{p,j}$ in the
column block of agent $j$, whereas $S_i$ and $P_i$ here are world-frame
projectors. What the present note adds is that this choice is forced rather
than convenient. By Lemma~\ref{lem:obstruction} the same restriction cannot be
expressed in $\Bper_{\mathcal G}$, and by Proposition~\ref{prop:unique}
$\Bnode_{\mathcal G}$ is the only matrix compatible with the null-column
requirement that Theorem~2 of \cite{michieletto2021} uses.

\section{Numerical checks}

The reference below is the Jacobian of the bearing map along admissible
variations, obtained by central differences on the manifold with attitudes
updated as $R_i\mapsto\exp([\varepsilon\omega_i]_\times)R_i$. It uses neither
construction.

Over $432$ frameworks covering the five manifolds at $n=3,4,6,9$ and sixteen
heterogeneous mixes at four graph densities, both \eqref{eq:peredge} and
\eqref{eq:pernode} reproduce this Jacobian on $\mathcal A$ to a worst relative
error of $1.1\times10^{-9}$, as \eqref{eq:F} requires of both. The identity
$\Bnode(I-AA^\top)=0$ holds to $2.2\times10^{-16}$, and the corresponding
quantity for \eqref{eq:peredge} is of order one. Over $1200$ random
heterogeneous frameworks with $n\in\{4,5,6\}$ we found the following.

\begin{center}
\begin{tabular}{lcc}
\toprule
 & edge-indexed & node-indexed\\
\midrule
$\rank\mathbf B=6n-q_v-q_i$ fails
    & $504/1198$
    & $0/1198$\\
rank test disagrees with the Jacobian
    & $67/1198$ $(5.6\%)$
    & $0/1198$\\
\bottomrule
\end{tabular}
\end{center}

For homogeneous formations the two constructions agree bitwise, checked over
$25$ graphs in each manifold at $n=3,4,6,8,12$, and the ranks match the closed
forms $2n-3$, $3n-4$, $3n-4$, $4n-5$, $6n-7$. Reproducing the case study of
Section~VI-B of \cite{michieletto2021} with \eqref{eq:peredge} returns
$\rank\Bper_{\mathcal K}=13$ with six null columns, the values reported
there, which suggests that the difference lies in the construction rather than
in our implementation of it.

\begin{thebibliography}{9}
\bibitem{michieletto2021}
G.~Michieletto, A.~Cenedese, and D.~Zelazo,
``A Unified Dissertation on Bearing Rigidity Theory,''
\emph{IEEE Transactions on Control of Network Systems}, vol.~8, no.~4,
pp.~1624--1636, Dec.~2021.
\href{https://doi.org/10.1109/TCNS.2021.3077712}{doi:10.1109/TCNS.2021.3077712}
\bibitem{Pozzan2021}
B.~Pozzan, G.~Michieletto, A.~Cenedese, and D.~Zelazo,
``Heterogeneous Formation Control: a Bearing Rigidity Approach,''
\emph{2021 60th IEEE Conference on Decision and Control (CDC)},
pp.~6451--6456, 2021.
\href{https://doi.org/10.1109/CDC45484.2021.9683374}{doi:10.1109/CDC45484.2021.9683374}
\end{thebibliography}
\end{document}
